% ===== 原文 pp.103–106（PDF p.113–116） =====
% [原文 p.103]
其中最后一行是 Vafa-Witten 作用量中 $\theta$ 角置零时瞬子项的贡献，我们已经略去了费米子项。截断作用量中不含场 $A_1$ 与 $A_2$ 的任何（沿 $\Sigma$ 求导的）动能项；也就是说，在 $\Sigma$ 上的有效作用量里它们可以被当作辅助场，并通过其运动方程积掉。这会把它们设成某些费米子模的组合，但既然我们一直略去费米子贡献，就可以简单地把 $A_1$ 与 $A_2$ 从有效作用量中丢弃。
% [Tutor Note: p.103 “不含 $A_1,A_2$ 的动能项”可按 p.101 的 $\epsilon$ 幂次计数核对：$F_{12}$ 不含任何 $C$ 方向指标，只带 $\epsilon^{+1}$，属被丢弃的次领头项，故 (8.3) 中没有 $F_{12}^2$。另注意 (8.3) 中 $2F_{13}F_{24}-2F_{14}F_{23}$ 正是“补全平方”的交叉项：$(F_{13}-F_{24})^2+(F_{14}+F_{23})^2+2F_{13}F_{24}-2F_{14}F_{23}=F_{13}^2+F_{14}^2+F_{23}^2+F_{24}^2$，可据此逐项验算]

\section{以 \texorpdfstring{$\cM_H(G,C)$}{MH(G,C)} 为靶空间的修正 A-模型}

现在我们把已经得到的场含量与有效作用量解释为 $\Sigma$ 上修正 A-模型的场含量与有效作用量。方便的做法是在施加 Hitchin 方程并按规范变换取商之前，先使用对 $(A_C,\varphi)$ 构成的空间；这是一个无限维的平坦簇
\[
\cW=\cA(G,C)\times\Omega^1(C,\mathrm{ad}\,P_C),
\]
其中 $\cA(G,C)$ 是 $P_C$ 上联络构成的空间。设 $z=\sigma^1+\ii\sigma^2$、$w=\sigma^3+\ii\sigma^4$ 分别是 $\Sigma$ 与 $C$ 上的复坐标。采用文献 [1] §4.1 的记号，我们取 $\cW$ 上的平坦度规为
\[
\dd s^2=-\frac{1}{4\pi}\int_C(\ii\,\dd w\wedge\dd\bar w)\;\Tr\left\{\delta A_w\otimes\delta A_{\bar w}+\delta A_{\bar w}\otimes\delta A_w+\delta\varphi_w\otimes\delta\varphi_{\bar w}+\delta\varphi_{\bar w}\otimes\delta\varphi_w\right\},
\]
其中 $\delta$ 表示场空间上的外微分。该度规关于 $\cW$ 上三个满足四元数关系的独立复结构 $I$、$J$、$K$ 是超 Kähler 的。我们用 $\omega_I$、$\omega_J$、$\omega_K$ 记相应的三个 Kähler 形式。（复结构与 Kähler 形式的显式表达式见文献 [1] §4.1；下面我们只写出 $\omega_I$。）

Hitchin 模空间也是超 Kähler 的。这一事实来自它作为 $\cW$ 关于丛 $P_C$ 的规范变换群所作的超 Kähler 商的描述，具体如下。考虑由元素 $\epsilon\in\Omega^0(C,\mathrm{ad}\,P_C)$ 生成的无穷小规范变换，其中我们把无穷维规范变换群的李代数等同于 $\mathrm{ad}\,P_C$ 的截面构成的向量空间。与辛形式 $\omega_I$、$\omega_J$、$\omega_K$ 对应的三个矩映射由下式给出
\begin{gather}
\mu_I(\epsilon)=\frac{1}{2\pi}\int_C(\ii\,\dd w\wedge\dd\bar w)\;\Tr\epsilon\left\{F-\tfrac12[\varphi,\varphi]\right\} \notag\\
\mu_J(\epsilon)=\frac{1}{2\pi}\int_C(\ii\,\dd w\wedge\dd\bar w)\;\Tr\epsilon\left\{D_w\varphi_{\bar w}+D_{\bar w}\varphi_w\right\} \notag\\
\mu_K(\epsilon)=\frac{\ii}{2\pi}\int_C(\ii\,\dd w\wedge\dd\bar w)\;\Tr\epsilon\left\{D_w\varphi_{\bar w}-D_{\bar w}\varphi_w\right\} \notag
\end{gather}
% [原文 p.104]
其中我们把矩映射取在 Lie 代数元素 $\epsilon$ 上求值。例如有
\[
\delta\mu_I(\epsilon)=-\iota_{V(\epsilon)}\,\omega_I,
\]
其中 $V(\epsilon)$ 是 $\cW$ 上对应于由 $\epsilon$ 生成的规范变换作用的向量场。Hitchin 模空间 $\cM_H(G,C)$ 正是水平集 $\mu_I(\epsilon)=\mu_J(\epsilon)=\mu_K(\epsilon)=0$（对一切 $\epsilon$）按规范变换群取的商，因此它本身就是一个超 Kähler 流形。在不可约点处（按假设，这里我们只考虑这些区域），$\cM(G,C)$ 的维数与其虚拟维数一致，后者定义为把 Hitchin 方程在该点处线性化所得某个椭圆微分复形的指数。文献 [64] 证明，该维数与规范群的维数及 $C$ 的亏格有如下关系
\[
\dim\cM_H(G,C)=4(g_C-1)\dim G .
\]
% [Tutor Note: p.104 的 $\dim\cM_H=4(g_C-1)\dim G$ 可用 $G=SU(2)$ 核对：$\dim G=3$ 给出 $\dim_{\R}\cM_H=12(g-1)$，而 Hitchin 模空间的复维数对 $SU(2)$ 为 $6(g-1)$（稳定丛 $3(g-1)$ 加 $\varphi\in H^0(\mathrm{ad}\,P\otimes K)$ 的 $3(g-1)$），实维数恰为两倍，一致。另注：此式中 $\dim G$ 指实维数]

与 Vafa-Witten 理论紧化相关的复结构是我们称为 $I$ 的那个，其相应的 Kähler 形式为
\[
\omega_I=-\frac{\ii}{2\pi}\int_C(\ii\,\dd w\wedge\dd\bar w)\;\Tr\left\{\delta A_{\bar w}\wedge\delta A_w-\delta\varphi_{\bar w}\wedge\delta\varphi_w\right\}.
\]
我们再次强调，上式中的 $\delta A$ 被视为场空间 $\cW$ 上的 1-形式，而 $\wedge$ 同时表示 $C$ 上与 $\cW$ 上的楔积。注意该 Kähler 形式被前面定义的 $U(1)$ 作用 $\cU_1:\cW\to\cW$ 保持，而且该作用关于整体定义的矩映射
\[
H=-\frac{1}{2\pi}\int_C(\ii\,\dd w\wedge\dd\bar w)\;\Tr\left\{\varphi_w\varphi_{\bar w}\right\}
\]
是哈密顿作用，此式视为 $\cW$ 上的实函数。容易验证 $\cU_1$ 把 $\mu$ 的零点集映入自身，因而下降为 $\cM_H(G,C)$ 上的哈密顿 $U(1)$ 作用。
% [Tutor Note: p.104 “容易验证”处补出论证：$\cU_1$ 保持 $A$ 不动而令 $\varphi_w\to\ee^{\ii\alpha}\varphi_w$、$\varphi_{\bar w}\to\ee^{-\ii\alpha}\varphi_{\bar w}$；于是 $[\varphi,\varphi]=2[\varphi_w,\varphi_{\bar w}]$ 不变，$D_{\bar w}\varphi_w=0$ 与 $D_w\varphi_{\bar w}=0$ 两式各自只乘一个相因子仍为零，故 $\mu_I=\mu_J=\mu_K=0$ 的零点集被保持；同时 $\omega_I$ 中 $\delta\varphi_{\bar w}\wedge\delta\varphi_w$ 的相因子相消，故 $\omega_I$ 不变]

我们来写下靶空间为 $\cM_H(G,C)$、Kähler 形式为 $\omega_I$、哈密顿 $U(1)$ 作用为 $\cU_1$ 的修正 A-模型的作用量。用靶空间上的复坐标 $X^i$ 与 $X^{\bar j}$，可把作用量 (7.7) 写成
\[
S=t\int_\Sigma(\ii\,\dd z\wedge\dd\bar z)\left\{2g_{i\bar j}\cD_{\bar z}X^i\cD_zX^{\bar j}+\text{fermions}\right\}+t\int_\Sigma\Phi^*(\omega_I)_{eq},
\]
其中 $t$ 是该二维理论的耦合常数，而 $\Phi^*(\omega_I)_{eq}$ 是等变辛形式在（相伴纤维丛的截面 $\Phi$ 下的）拉回
\[
(\omega_I)_{eq}=\omega_I-\zeta H\;\in\;\Omega^2_{U(1)}(\cM_H).
\]
% [原文 p.105]
其中 $\zeta$ 是 $\fu(1)$ 的二次生成元，在拉回下它被替换为自旋联络的曲率。（关于相伴纤维丛的截面下等变形式之拉回这一概念，详见附录 C。）协变导数为
\begin{gather}
\cD_{\bar z}X^i=\partial_{\bar z}X^i+\omega_{\bar z}V^i=\partial_{\bar z}X^i-\ii(\partial_{\bar z}f_\Sigma)V^i \notag\\
\cD_zX^{\bar j}=\partial_zX^{\bar j}+\omega_zV^{\bar j}=\partial_zX^{\bar j}+\ii(\partial_zf_\Sigma)V^{\bar j} \notag
\end{gather}
其中我们代入了 (7.6) 以给出 $\Sigma$ 上自旋联络的分量 $\omega_{\bar z}$ 与 $\omega_z$，而 $V$ 是生成 $U(1)$ 作用的向量场。

在复结构 $I$ 下，$\cW$ 上的线性、全纯与反全纯函数由下式给出
\begin{gather}
X^i\;\longrightarrow\;(A_{\bar w},\varphi_w)_c \notag\\
X^{\bar j}\;\longrightarrow\;(A_w,\varphi_{\bar w})_c \notag
\end{gather}
其中 $c\in C$ 是求值点。协变导数为
\begin{gather}
\cD_{\bar z}X^i\;\longrightarrow\;\left(\partial_{\bar z}A_{\bar w},\;\partial_{\bar z}\varphi_w+(\partial_{\bar z}f_\Sigma)\varphi_w\right)_c \notag\\
\cD_zX^{\bar j}\;\longrightarrow\;\left(\partial_zA_w,\;\partial_z\varphi_{\bar w}+(\partial_zf_\Sigma)\varphi_{\bar w}\right)_c . \notag
\end{gather}

把这些表达式（以及上面写出的度规）代入我们对 $S$ 的表达式，可得
\begin{equation*}
\begin{aligned}
S=-\frac{t}{2\pi}\int_\Sigma(\ii\,\dd z\wedge\dd\bar z)\int_C(\ii\,\dd w\wedge\dd\bar w)\;\Tr\Big\{&2\partial_zA_w\partial_{\bar z}A_{\bar w}+2\left(\partial_z\varphi_{\bar w}+(\partial_zf_\Sigma)\varphi_{\bar w}\right)\left(\partial_{\bar z}\varphi_w+(\partial_{\bar z}f_\Sigma)\varphi_w\right)\\
&+\text{fermions}\Big\}+t\int_\Sigma\Phi^*(\omega_I)_{eq}.
\end{aligned}
\end{equation*}

为了与上一节得到的有效作用量建立联系，我们对等变拓扑项稍作“加工”，利用
\[
(\omega_I)_{eq}=\omega'_I+\delta_C\lambda_I,
\]
其中 2-形式 $\omega'_I$ 与 1-形式 $\lambda_I$ 由下式给出
\begin{gather}
\omega'_I=-\frac{\ii}{2\pi}\int_C(\ii\,\dd w\wedge\dd\bar w)\;\Tr\left\{\delta A_{\bar w}\wedge\delta A_w\right\}, \notag\\
\lambda_I=-\frac{\ii}{4\pi}\int_C(\ii\,\dd w\wedge\dd\bar w)\;\Tr\left\{\varphi_w\,\delta\varphi_{\bar w}-\varphi_{\bar w}\,\delta\varphi_w\right\} \notag
\end{gather}
% [原文 p.106]
而 $\delta_C=\delta-\zeta\iota_V$ 是 $\cW$ 上的 Cartan 等变微分。也就是说，等变辛形式 $(\omega_I)_{eq}$ 与 2-形式 $\omega'_I$ 等变上同调，由 §C.2 的一个定理，在截面 $\Phi$ 下拉回并在 $\Sigma$ 上积分所得的数值将完全相同：
\[
\int_\Sigma\Phi^*(\omega_I)_{eq}=\int_\Sigma\Phi^*(\omega'_I).
\]
% [Tutor Note: p.105–106 的等式 $(\omega_I)_{eq}=\omega'_I+\delta_C\lambda_I$ 可逐项验算：$\delta\lambda_I=-\frac{\ii}{4\pi}\int\Tr\{\delta\varphi_w\wedge\delta\varphi_{\bar w}-\delta\varphi_{\bar w}\wedge\delta\varphi_w\}=\frac{\ii}{2\pi}\int\Tr\{\delta\varphi_{\bar w}\wedge\delta\varphi_w\}=\omega_I-\omega'_I$；又 $\iota_V\lambda_I=H$（用 $V(\varphi_w)=\ii\varphi_w$、$V(\varphi_{\bar w})=-\ii\varphi_{\bar w}$ 直接代入），故 $\omega'_I+\delta_C\lambda_I=\omega'_I+\delta\lambda_I-\zeta\iota_V\lambda_I=\omega_I-\zeta H$，与 $(\omega_I)_{eq}$ 的定义一致。这是原文“we massage … slightly”处省略的全部计算]

作了这个替换之后，我们最终得到以 $\cM_H(G,C)$ 为靶空间的修正 A-模型的作用量：
\begin{equation*}
\begin{aligned}
S=-\frac{t}{2\pi}\int_\Sigma(\ii\,\dd z\wedge\dd\bar z)\int_C(\ii\,\dd w\wedge\dd\bar w)\;\Tr\Big\{&\partial_zA_w\partial_{\bar z}A_{\bar w}+\partial_{\bar z}A_w\partial_zA_{\bar w}\\
&+2\left(\partial_z\varphi_{\bar w}+(\partial_zf_\Sigma)\varphi_{\bar w}\right)\left(\partial_{\bar z}\varphi_w+(\partial_{\bar z}f_\Sigma)\varphi_w\right)+\text{fermions}\Big\}.
\end{aligned}
\end{equation*}
若把二维耦合常数 $t$ 认定为
\[
t=\frac{4\pi}{e^2},
\]
则上述作用量恰好就是我们在 (8.3) 中为截断的 Vafa-Witten 作用量写下的东西。有效二维 TFT 定域在满足下式的截面 $\Phi$（即 $(A_C,\varphi)$ 的构形）上：
\[
\partial_{\bar z}A_{\bar w}=\partial_zA_w=\partial_{\bar z}\varphi_w+(\partial_{\bar z}f_\Sigma)\varphi_w=\partial_z\varphi_{\bar w}+(\partial_zf_\Sigma)\varphi_{\bar w}=0 .
\]
% [Tutor Note: 把 (8.3) 与 p.116 的 $S$ 逐项对照时，测度相差因子 4（$\ii\dd z\wedge\dd\bar z=2\,\dd\sigma^1\dd\sigma^2$，故 $\int_\Sigma\int_C=4\int_M\dd^4\sigma$），再叠加 $\Tr$ 的归一化（作者在 [1] §4.1 中的约定），字面系数会与 $t=4\pi/e^2$ 差一个常数；这不影响任何物理结论，核对时以定域化方程（p.106 末四式）与 Hitchin 约束为准]

\para{膜（Branes）}

当 $\Sigma$ 有非平凡边界时，必须沿边界的每个连通分支对场施加边界条件。最简单的边界条件对应于靶空间的拉格朗日子流形 $Y$，并且还要求它被哈密顿 $U(1)$ 作用保持。\footnote{另外还要求嵌入子流形的等变 Maslov 指标为零，见 §7.5。}（细节见 §7.5。）如文献 [35] 所讨论，Hitchin 模空间 $\cM_H(G,C)$ 关于辛形式 $\omega_I$ 有几类简单的拉格朗日子流形。然而这些例子中没有一个关于 $\cU_1$ 对称性的作用是保持不变的。
